\scanpage{26}
\mspage{-2.22-}
% archive sheet number: 25

of some torus $T\subset H$. If $\epsilon_1\neq\epsilon_2$, then there do not exist positive integers $n$ and $m$ such that $\lambda_1(\handwritten{\alpha^n})=\lambda_2(\handwritten{\alpha^m})$, all $\alpha\in\handwritten{\mathbb{G}}_m$. But then there is a character $\handwritten{\chi}$ of $T$ such that $\handwritten{\chi}\circ\lambda_1=0$, (the trivial homomorphism from $\handwritten{\mathbb{G}}_m$ to $\handwritten{\mathbb{G}}_m$), and $\handwritten{\chi}\circ\lambda_2\neq0$. Since the kernel of $\phi$ is finite, there is an integer $N$ such that $\handwritten{\chi}^N$ is induced by a character $\handwritten{\chi'}$ of $\phi(T)$. Now suppose that $\phi_*\epsilon_1=\phi_*\epsilon_2$. Then $P(\phi\circ\lambda_1)=P(\phi_*\epsilon_1)=P(\phi_*\epsilon_2)=P(\phi\circ\lambda_2)$. Let $T^{\handwritten{\prime\prime}}$ be the intersection of $\phi(T)$, with the radical of $P(\phi\circ\lambda_1)$, and let $\handwritten{\chi''}$ be the restriction of $\handwritten{\chi'}$ to $T$. According to the structure theory of parabolic subgroups (cf.\hspace{4em}), some power of $\handwritten{\chi''}$ extends to a character $\handwritten{\chi'''}$ of $P(\phi\circ\lambda_1)$. But since $\phi\circ\lambda_1$ and $\phi\circ\lambda_2$ are 1-parameter subgroups of $P(\phi\circ\lambda_1)$, hence of $T^{\handwritten{\prime\prime}}$, we have $\handwritten{\chi'''}\circ(\phi\circ\lambda_1)=0$, and $\handwritten{\chi'''}\circ(\phi\circ\lambda_2)\neq0$. But then, referring to the definition of $\Delta(G)$, one sees that this precludes the possibility $\Delta(\phi\circ\lambda_1)=\Delta(\phi\circ\lambda_2)$. Therefore $\phi_*\epsilon_1\neq\phi_*\epsilon_2$. \uline{QED}

\uline{Corollary} \uline{2.11}: If $\phi$ is an isogeny, $\phi_*$ is an isomorphism.

\uline{Proof}: It remains to show that $\phi_*$ is surjective. But if $\lambda$ is any 1-parameter subgroup of $G$, then for some positive integer $n$, and for some 1-parameter subgroup $\mu$ of $H$, $\phi\circ\mu(\alpha)=\lambda(\alpha^{\handwritten{\unsure{h|n}}})$, $\alpha\in\handwritten{\mathbb{G}}_m$. \uline{QED}

Suppose we apply Proposition 2.1\deleted{\unsure{0|1}}\handwritten{\unsure{1|0}} to the inclusion of a maximal torus $T$ in $G$. This induces an inclusion of $\Delta(T)$ in $\Delta(G)$. But $\Delta(T)$ is very simple - it is essentially the set of rational points on a sphere. In fact, it follows immediately from the definition that $\Delta(T)$ is the set of rays in $\unsure{\Gamma|\Pi}^{\handwritten{\mathbb{R}}}(T)$ which contain a point of $\Gamma^{\handwritten{\mathbb{Q}}}(T)$.
